\documentclass[11pt]{amsart}
\usepackage{amsmath,amssymb,amsthm}
\usepackage{booktabs}
\usepackage{xcolor}
\usepackage[no-math]{fontspec}
\setmonofont{DejaVuSansMono.ttf}[Scale=0.82]
\usepackage[margin=1.1in]{geometry}
\usepackage[colorlinks=true,linkcolor=blue!50!black,citecolor=blue!50!black,urlcolor=blue!50!black]{hyperref}
\usepackage{fancyvrb}

\newtheorem{theorem}{Theorem}[section]
\newtheorem{proposition}[theorem]{Proposition}
\newtheorem{lemma}[theorem]{Lemma}
\newtheorem{corollary}[theorem]{Corollary}
\newtheorem{conjecture}[theorem]{Conjecture}
\theoremstyle{definition}
\newtheorem{definition}[theorem]{Definition}
\theoremstyle{remark}
\newtheorem{remark}[theorem]{Remark}

\newcommand{\lf}{\lambda}
\newcommand{\Sh}{\mathcal{S}}
\newcommand{\Prn}{\mathcal{P}}
\newcommand{\cstar}{c^{*}}
\newcommand{\join}[2]{[#1,#2]}
\newcommand{\binomtwo}[1]{\tbinom{#1}{2}}
\DeclareMathOperator{\shape}{sh}

\title[Non-equivalent ancestral configurations]{The maximum number of non-equivalent ancestral configurations:\\ a proof of a conjecture of Disanto and Rosenberg}

\author{Claude Opus 5.5}
\author{GPT-6-Astra}
\author{Claude Fable 5.1}
\thanks{The authors are AI models: Claude Opus~5.5 and Claude Fable~5.1 (Anthropic) and GPT-6-Astra (OpenAI). Source files and Lean code: \url{https://github.com/Kevin-OS7/math}.}

\date{October 29, 2026}

\begin{document}

\begin{abstract}
Ancestral configurations are the combinatorial objects that underlie the computation of gene tree probabilities under the multispecies coalescent.
For a species tree with $n$ leaves and a gene tree of the same labeled topology, Disanto and Rosenberg (Bull.\ Math.\ Biol.\ 2019) studied the number $\cstar(T)$ of \emph{non-equivalent} root ancestral configurations, i.e.\ root configurations counted up to isomorphism of the tree shapes they induce.
They conjectured that the maximum of $\cstar(T)$ over all tree shapes $T$ with $n\ge 2$ leaves equals
\[
\Bigl\lfloor \tfrac{1}{27}\,3^{\lfloor n/3\rfloor}\bigl(2r^2+8r+23\bigr)\Bigr\rfloor,\qquad r=n-3\lfloor n/3\rfloor,
\]
and verified this for $n\le 20$.
We prove the conjecture for all $n$.
We further show that for every $n\ge 9$ the maximizing tree shape is unique, and identify it; that the exponential growth constant of the maximum is exactly $3^{1/3}$, which settles the value of a constant left undetermined by Disanto and Rosenberg; and that the minimum of $\cstar(T)$ is $n-1$, attained by exactly two tree shapes for every $n\ge 4$.
The proof rests on an exact product formula for the number of pruned shapes of a tree together with an elementary but delicate arithmetic inequality.
The main theorem, the uniqueness statement and the lower bound $n-1$ have been formally verified in the Lean~4 proof assistant with the Mathlib library.
\end{abstract}

\maketitle

\section{Introduction}

Under the multispecies coalescent model, the probability of a gene tree topology given a species tree can be computed by dynamic programming over \emph{ancestral configurations}: the sets of gene lineages that can be present at the nodes of the species tree \cite{Wu2012}.
The number of ancestral configurations therefore governs the cost of such computations, and its combinatorics has been studied in a series of papers by Disanto, Rosenberg and coauthors \cite{DR2017,DR2019,DFPR2022,DFHPR2024,LR2024}.

For a \emph{matching} gene tree and species tree (the two have the same labeled topology), Disanto and Rosenberg \cite{DR2017} enumerated the root ancestral configurations and showed that their maximum number over trees with $n$ leaves grows like $k_0^{\,n}$ with $k_0\approx 1.5028$.
In \cite{DR2019} they considered a coarser count suggested by Wu's algorithm \cite{Wu2012}: two configurations at a node $\kappa$ are \emph{equivalent} if the tree shapes obtained by cutting the subtree at $\kappa$ along the two configurations are isomorphic.
Writing $\cstar(T)$ for the number of non-equivalent root configurations when the species tree has unlabeled shape $T$, they exhibited a family of shapes $u_n$ with
\[
\cstar(u_n)=\Bigl\lfloor 3^{\lfloor n/3\rfloor}\tfrac{2r^2+8r+23}{27}\Bigr\rfloor,\qquad r=n-3\lfloor n/3\rfloor,
\]
found by exhaustive computation that $u_n$ attains the maximum for $2\le n\le 20$, and stated the following conjecture \cite[Sect.~4.3, after eq.~(16)]{DR2019}.

\begin{conjecture}[Disanto--Rosenberg]\label{conj:DR}
For every $n\ge 2$,
\[
M_n:=\max_{|T|=n}\cstar(T)=f(n):=\Bigl\lfloor \frac{3^{\lfloor n/3\rfloor}\,(2r^2+8r+23)}{27}\Bigr\rfloor,\qquad r=n\bmod 3 .
\]
\end{conjecture}

They also proved $3^{n/3}\ll M_n\ll k_0^{\,n}$, so that the growth constant $k_1^*=\lim M_n^{1/n}$ (if it exists) satisfies $3^{1/3}\le k_1^*\le k_0\approx 1.5028$, and remarked that they ``did not prove a general result describing the unlabeled shapes of trees that give rise to the most non-equivalent root configurations'' \cite[Sect.~6]{DR2019}.

\subsection*{Results}
We prove Conjecture~\ref{conj:DR} and describe the extremal trees.
Let $\lf$ denote the one-leaf tree, $[x,y]$ the (unordered) tree whose root has subtrees $x$ and $y$, $\kappa_2=[\lf,\lf]$ and $\kappa_3=[\lf,\kappa_2]$.
Define shapes $u_n$ by $u_1=\lf$, $u_2=\kappa_2$, $u_3=\kappa_3$, $u_4=[\lf,\kappa_3]$ and, for $n\ge 5$,
\[
u_n=\begin{cases}[\kappa_2,u_{n-2}] & n\in\{5,6,7,8,9,11\},\\ [\kappa_3,u_{n-3}] & n=10 \text{ or } n\ge 12.\end{cases}
\]
(For $n\ge 9$ these are the trees of \cite[Sect.~3.2 and eq.~(14)]{DR2019}.)

\begin{theorem}\label{thm:main}
For every $n\ge 2$ we have $M_n=f(n)$, and $\cstar(u_n)=f(n)$.
\end{theorem}

\begin{theorem}\label{thm:unique}
For every $n\ge 9$, $u_n$ is the only tree shape with $n$ leaves satisfying $\cstar(T)=M_n$.
\end{theorem}

\begin{corollary}\label{cor:growth}
For $n\ge 9$, $M_n=C_r\,3^{(n-r)/3-3}$ with $(C_0,C_1,C_2)=(23,33,47)$ and $r=n\bmod 3$; in particular $M_{n+3}=3M_n$ for $n\ge 9$ and $\lim_{n\to\infty} M_n^{1/n}=3^{1/3}$.
Thus the growth constant $k_1^*$ of \cite{DR2019} equals $3^{1/3}\approx 1.4422$.
\end{corollary}

\begin{proposition}\label{prop:min}
For every tree shape $T$ with $n$ leaves, $\cstar(T)\ge n-1$.
Equality holds for exactly one shape when $n\le 3$ and for exactly two shapes when $n\ge 4$: the caterpillar, and the tree obtained from the caterpillar with $n-2$ leaves by replacing its cherry with the balanced tree $[\kappa_2,\kappa_2]$.
\end{proposition}

Proposition~\ref{prop:min} complements the observation in \cite[Sect.~3.2]{DR2019} that, for $n\le 20$, the caterpillar minimizes $\cstar$: the minimum value is $n-1$ for all $n$, and for $n\ge 4$ the caterpillar is not the only minimizer.

\subsection*{Method}
The key observation (Section~\ref{sec:prunings}) is that $\cstar(T)+1$ is the number $s(T)$ of distinct shapes of \emph{prunings} of $T$, and that for $T=[L,R]$
\begin{equation}\label{eq:product}
s(T)=1+s(L)\,s(R)-\binomtwo{k},\qquad k=\bigl|\Sh(L)\cap\Sh(R)\bigr|,
\end{equation}
where $\Sh(\cdot)$ denotes the set of pruned shapes.
Every tree with at least three leaves has $\lf,\kappa_2,\kappa_3$ among its pruned shapes, so $k\ge\min(|L|,|R|,3)$.
Theorem~\ref{thm:main} then reduces to the arithmetic inequality of Lemma~\ref{lem:arith},
\[
\bigl(F(a)+1\bigr)\bigl(F(b)+1\bigr)-\binomtwo{\min(a,b,3)}\le F(a+b)\qquad(a,b\ge 1),
\]
where $F(1)=0$ and $F(n)=f(n)$ for $n\ge 2$.
Formula~\eqref{eq:product} generalizes \cite[eq.~(9)]{DR2019}, which covers the case where one subtree is displayed in the other.

\subsection*{Formal verification}
Theorems~\ref{thm:main} and \ref{thm:unique} and the inequality $\cstar(T)\ge n-1$ of Proposition~\ref{prop:min} have been formalized and machine-checked in Lean~4 \cite{Lean4} using Mathlib \cite{Mathlib}; see Section~\ref{sec:lean}.

\section{Trees, cuts and prunings}\label{sec:prunings}

All trees are rooted binary trees in which every internal node has exactly two children; $|T|$ denotes the number of leaves.
Two trees are isomorphic if they are equal up to exchanging the two children at some internal nodes; a \emph{shape} is an isomorphism class.
For shapes $x,y$ we write $[x,y]=[y,x]$ for the shape whose root subtrees have shapes $x$ and $y$.
A shape with at least two leaves determines the unordered pair (multiset) $\{x,y\}$ of its root subtrees.

\begin{definition}
A \emph{cut} of $T$ is a set $X$ of nodes of $T$ such that every path from the root $\rho$ to a leaf contains exactly one node of $X$.
The \emph{pruning} $T|_X$ is obtained from $T$ by deleting all nodes strictly below the nodes of $X$ (so the nodes of $X$ become the leaves of $T|_X$).
We write $\Prn(T)$ for the set of prunings of $T$, $\Sh(T)$ for the set of their shapes, and $s(T)=|\Sh(T)|$.
\end{definition}

\begin{lemma}[Reduction]\label{lem:reduction}
Let the species tree and the matching gene tree have shape $T$.
Then the number of non-equivalent root ancestral configurations in the sense of \cite{DR2019} is $\cstar(T)=s(T)-1$.
\end{lemma}

\begin{proof}
For a matching pair, the gene lineages present at a node $\kappa$ of the species tree are identified with nodes of the gene tree $=$ species tree, and the root configurations are precisely the sets of gene lineages entering the root from its two child branches.
From the left child branch either the single lineage of the left child $\rho_\ell$ enters (all coalescences in the left clade have occurred), or a configuration at $\rho_\ell$ enters; by induction, the configurations at a node $\kappa$ are exactly the cuts of the subtree $T_\kappa$ different from $\{\kappa\}$.
This is the content of the recursion $c(\rho)=[c(\rho_\ell)+1][c(\rho_r)+1]$, $c(\text{leaf})=0$, of \cite[eq.~(5)]{DR2019} (see also \cite[Sect.~2]{DR2017}).
Hence root configurations correspond bijectively to the cuts $X\neq\{\rho\}$ of $T$.
By definition \cite[Sect.~3.1]{DR2019}, two configurations are equivalent if and only if the restricted shapes obtained by ``removing from $T$ all nodes that remain strictly below'' the configuration are isomorphic, i.e.\ if and only if $T|_X$ and $T|_{X'}$ have the same shape.
The cut $\{\rho\}$ yields $\lf$, and every other cut has at least two elements, so its pruning is not $\lf$.
Therefore $\cstar(T)=|\Sh(T)\setminus\{\lf\}|=s(T)-1$.
\end{proof}

The cuts of $\lf$ are $\{\lf\}$; the cuts of $T=[L,R]$ are $\{\rho\}$ and the sets $X_L\cup X_R$ with $X_L,X_R$ cuts of $L,R$; and $T|_{X_L\cup X_R}=[L|_{X_L},R|_{X_R}]$.
Hence
\begin{equation}\label{eq:recursion}
\Sh(\lf)=\{\lf\},\qquad \Sh([L,R])=\{\lf\}\cup\bigl\{[x,y]: x\in\Sh(L),\ y\in\Sh(R)\bigr\}.
\end{equation}

\begin{lemma}[Product formula]\label{lem:product}
For $T=[L,R]$ and $k=|\Sh(L)\cap\Sh(R)|$ we have $s(T)=1+s(L)s(R)-\binomtwo{k}$.
\end{lemma}

\begin{proof}
By \eqref{eq:recursion}, $s(T)-1$ is the number of distinct shapes $[x,y]$ with $(x,y)\in\Sh(L)\times\Sh(R)$, none of which is $\lf$.
Since $[x,y]$ determines the multiset $\{x,y\}$, two pairs $(x,y)\neq(x',y')$ give the same shape if and only if $x'=y$, $y'=x$ and $x\neq y$; in that case both $x$ and $y$ lie in $\Sh(L)\cap\Sh(R)$.
Conversely, every $2$-element subset $\{x,y\}\subseteq\Sh(L)\cap\Sh(R)$ arises from exactly the two pairs $(x,y)$ and $(y,x)$.
So the map $(x,y)\mapsto[x,y]$ is two-to-one over the $\binomtwo{k}$ such subsets and one-to-one elsewhere, giving $s(T)-1=s(L)s(R)-\binomtwo{k}$.
\end{proof}

\begin{lemma}[Small prunings]\label{lem:small}
If $|T|=m$, then for every $1\le j\le m$ the set $\Sh(T)$ contains a shape with exactly $j$ leaves.
The shapes with $1,2,3$ leaves are exactly $\lf,\kappa_2,\kappa_3$.
Consequently, for $T=[L,R]$ we have $|\Sh(L)\cap\Sh(R)|\ge\min(|L|,|R|,3)$.
\end{lemma}

\begin{proof}
The set of leaves is a cut with $T|_X=T$.
If $T|_X$ has $j\ge 2$ leaves, it contains two sibling leaves $u,v$ with parent $p$; then $(X\setminus\{u,v\})\cup\{p\}$ is a cut whose pruning has $j-1$ leaves.
The second statement is clear, and the third follows from the first two.
\end{proof}

\section{An arithmetic inequality}\label{sec:arith}

Let $F(1)=0$ and $F(n)=f(n)$ for $n\ge 2$ (note that the formula for $f$ would give $f(1)=1$).
The first values are
\[
\bigl(F(1),\dots,F(11)\bigr)=(0,1,2,3,5,7,11,15,23,33,47).
\]
For $n=3m+r\ge 9$ ($0\le r\le 2$) the floor is not needed:
\begin{equation}\label{eq:closed}
F(n)=C_r\,3^{m-3},\qquad (C_0,C_1,C_2)=(23,33,47).
\end{equation}
In particular $F(n)=3F(n-3)$ for $n\ge 12$, and one checks $F(10)=3F(7)$ and $F(n)=2F(n-2)+1$ for $n\in\{5,6,7,8,9,11\}$.

\begin{lemma}\label{lem:arith}
For all integers $a,b\ge 1$,
\begin{equation}\label{eq:arith}
\bigl(F(a)+1\bigr)\bigl(F(b)+1\bigr)-\binomtwo{\min(a,b,3)}\ \le\ F(a+b).
\end{equation}
Moreover, if $a\le b$ and $a+b\ge 15$, equality holds if and only if $a=3$.
\end{lemma}

\begin{proof}
By symmetry assume $a\le b$.

\emph{Case 1: $a\ge 9$.}
Write $a=3m_a+r_a$, $b=3m_b+r_b$, $X=3^{m_a-3}\ge1$, $Y=3^{m_b-3}\ge 1$.
By \eqref{eq:closed} the left side of \eqref{eq:arith} equals $C_{r_a}C_{r_b}XY+C_{r_a}X+C_{r_b}Y-2$, while the right side equals $27\cdot 3^{c}\,C_{r'}XY$, where $r'\equiv r_a+r_b\pmod 3$, $0\le r'\le2$, and $c=\lfloor (r_a+r_b)/3\rfloor$.
The difference $D=27\cdot 3^{c}C_{r'}-C_{r_a}C_{r_b}$ equals
\[
92,\ 132,\ 188,\ 180,\ 312,\ 464\quad\text{for}\quad (r_a,r_b)=(0,0),(0,1),(0,2),(1,1),(1,2),(2,2),
\]
which exceeds $C_{r_a}+C_{r_b}\in\{46,56,70,66,80,94\}$ respectively.
Since $C_{r_a}X+C_{r_b}Y\le(C_{r_a}+C_{r_b})XY<D\,XY$, the inequality \eqref{eq:arith} is strict.

\emph{Case 2: $a\le 8<9\le b$.}
Write $b=3m+r$ and $Y=3^{m-3}\ge 1$.
Then $a+b\ge 10$ and, by \eqref{eq:closed}, both sides of \eqref{eq:arith} are affine in $Y$: the left side is $\alpha Y+\beta$ and the right side is $\gamma Y$, where
\[
\alpha=(F(a)+1)C_r,\qquad \beta=F(a)+1-\binomtwo{\min(a,3)},\qquad \gamma=C_{r'}\,3^{\lfloor (a+r)/3\rfloor},\quad r'=(a+r)\bmod 3.
\]
Table~\ref{tab:case2} lists these numbers.
In all $24$ cases $\gamma-\alpha\ge\beta\ge0$, hence $\gamma Y-\alpha Y-\beta\ge(\gamma-\alpha)-\beta\ge 0$.
Equality forces $\gamma-\alpha=\beta$ and either $\beta=0$ or $Y=1$; by the table this happens only for $a=3$ (all $b$), or for $(a,b)\in\{(2,9),(5,9)\}$, which have $a+b<15$.

\emph{Case 3: $a\le b\le 8$.}
This is a finite check of $36$ cases; equality occurs exactly for $(a,b)\in\{(1,1),(1,2),(1,3),(2,2),(2,3),(2,4),(2,5),(2,6),(3,5),(2,7),(3,7),(5,5),(5,7),(7,7)\}$, none of which has $a+b\ge 15$ (for instance $(a,b)=(7,8)$ gives $189<207$ and $(8,8)$ gives $253<297$).
\end{proof}

\begin{table}[t]
\centering
\small
\begin{tabular}{ccc|ccc|ccc|ccc}
\toprule
& & & \multicolumn{3}{c|}{$r=0$} & \multicolumn{3}{c|}{$r=1$} & \multicolumn{3}{c}{$r=2$}\\
$a$ & $F(a)$ & $\beta$ & $\alpha$ & $\gamma$ & $\gamma-\alpha$ & $\alpha$ & $\gamma$ & $\gamma-\alpha$ & $\alpha$ & $\gamma$ & $\gamma-\alpha$\\
\midrule
1 & 0 & 1 & 23 & 33 & 10 & 33 & 47 & 14 & 47 & 69 & 22 \\
2 & 1 & 1 & 46 & 47 & 1 & 66 & 69 & 3 & 94 & 99 & 5 \\
3 & 2 & 0 & 69 & 69 & 0 & 99 & 99 & 0 & 141 & 141 & 0 \\
4 & 3 & 1 & 92 & 99 & 7 & 132 & 141 & 9 & 188 & 207 & 19 \\
5 & 5 & 3 & 138 & 141 & 3 & 198 & 207 & 9 & 282 & 297 & 15 \\
6 & 7 & 5 & 184 & 207 & 23 & 264 & 297 & 33 & 376 & 423 & 47 \\
7 & 11 & 9 & 276 & 297 & 21 & 396 & 423 & 27 & 564 & 621 & 57 \\
8 & 15 & 13 & 368 & 423 & 55 & 528 & 621 & 93 & 752 & 891 & 139 \\
\bottomrule
\end{tabular}
\medskip
\caption{The constants of Case~2 in the proof of Lemma~\ref{lem:arith} ($b=3m+r\ge 9$).}
\label{tab:case2}
\end{table}

\section{Proof of Theorem~\ref{thm:main}}\label{sec:main}

\begin{proposition}[Upper bound]\label{prop:upper}
Every tree $T$ satisfies $s(T)\le F(|T|)+1$.
Hence $\cstar(T)\le f(n)$ whenever $|T|=n\ge 2$.
\end{proposition}

\begin{proof}
Induction on $T$.
For $T=\lf$, $s(T)=1=F(1)+1$.
For $T=[L,R]$ with $a=|L|$, $b=|R|$, Lemmas~\ref{lem:product} and \ref{lem:small} and the induction hypothesis give
\[
s(T)=1+s(L)s(R)-\binomtwo{k}\le 1+(F(a)+1)(F(b)+1)-\binomtwo{\min(a,b,3)}\le 1+F(a+b),
\]
using that $k\mapsto\binomtwo{k}$ is nondecreasing and $k\ge\min(a,b,3)$, and Lemma~\ref{lem:arith} in the last step.
The second statement follows from Lemma~\ref{lem:reduction}.
\end{proof}

\begin{proposition}[Lower bound]\label{prop:lower}
$s(u_n)=F(n)+1$ for all $n\ge 1$; hence $\cstar(u_n)=f(n)$ for $n\ge 2$.
\end{proposition}

\begin{proof}
Directly, $s(u_1)=1$, $s(u_2)=2$, $s(u_3)=3$ and $s(u_4)=4$.
If $|X|\ge 2$, then $\Sh(\kappa_2)=\{\lf,\kappa_2\}\subseteq\Sh(X)$ by Lemma~\ref{lem:small}, so $k=2$ in Lemma~\ref{lem:product} and $s([\kappa_2,X])=1+2s(X)-1=2s(X)$.
If $|X|\ge 3$, then $\Sh(\kappa_3)=\{\lf,\kappa_2,\kappa_3\}\subseteq\Sh(X)$, so $k=3$ and $s([\kappa_3,X])=3s(X)-2$.
Since $F(n)+1=2(F(n-2)+1)$ for $n\in\{5,6,7,8,9,11\}$ and $F(n)+1=3(F(n-3)+1)-2$ for $n=10$ and $n\ge 12$, induction gives $s(u_n)=F(n)+1$.
\end{proof}

Propositions~\ref{prop:upper} and \ref{prop:lower} prove Theorem~\ref{thm:main}, and Corollary~\ref{cor:growth} follows from \eqref{eq:closed}.
The values $M_2,\dots,M_{20}$, namely 1, 2, 3, 5, 7, 11, 15, 23, 33, 47, 69, 99, 141, 207, 297, 423, 621, 891, 1269, agree with the exhaustive computation of \cite[Fig.~5]{DR2019}.

\section{Uniqueness of the maximizer}\label{sec:unique}

\begin{proof}[Proof of Theorem~\ref{thm:unique}]
For $9\le n\le 14$ we verified by exhaustive enumeration of all tree shapes ($46$, $98$, $207$, $451$, $983$ and $2179$ shapes, respectively) that $u_n$ is the unique maximizer; the enumeration (a short program computing $\Sh(T)$ via \eqref{eq:recursion} for every shape) reproduces the maxima of \cite[Fig.~5]{DR2019} for all $n\le 19$ and finds a unique maximizer for every $9\le n\le 19$.
(These base cases are also verified inside the Lean formalization; see Section~\ref{sec:lean}.)

Let $n\ge 15$ and let $T=[L,R]$ with $a=|L|\le b=|R|$ satisfy $\cstar(T)=f(n)$.
In the chain of inequalities in the proof of Proposition~\ref{prop:upper} all inequalities are then equalities.
Equality throughout forces $s(L)=F(a)+1$, $s(R)=F(b)+1$, $k=\min(a,b,3)$, and equality in Lemma~\ref{lem:arith}.
By Lemma~\ref{lem:arith}, $a=3$, so $L=\kappa_3$ (the only shape with three leaves).
In particular $s(R)=F(n-3)+1$, i.e.\ $R$ is a maximizer for $n-3\ge 12$; by induction $R=u_{n-3}$, hence $T=[\kappa_3,u_{n-3}]=u_n$.
\end{proof}

For $n\le 8$ the maximizers are not unique in general; the complete lists are
$n=4$: $[\lf,\kappa_3]$, $[\kappa_2,\kappa_2]$;
$n=6$: $[\kappa_2,[\lf,\kappa_3]]$, $[\kappa_2,[\kappa_2,\kappa_2]]$;
$n=8$: $[\kappa_2,[\kappa_2,[\lf,\kappa_3]]]$, $[\kappa_2,[\kappa_2,[\kappa_2,\kappa_2]]]$, $[\kappa_3,[\kappa_2,\kappa_3]]$;
and $u_n$ is the unique maximizer for $n\in\{2,3,5,7\}$.

\section{The minimum}\label{sec:min}

\begin{proof}[Proof of Proposition~\ref{prop:min}]
By Lemma~\ref{lem:small}, $\Sh(T)$ contains shapes with $1,2,\dots,n$ leaves, so $s(T)\ge n$ and $\cstar(T)\ge n-1$.
Call $T$ \emph{minimal} if $s(T)=|T|$, i.e.\ $\Sh(T)$ contains exactly one shape with $j$ leaves for each $j\le|T|$.
Let $T=[L,R]$ with $a=|L|\le b=|R|$.
If $a\ge 2$ and $b\ge 3$, then $[\kappa_2,\kappa_2]$ and $[\lf,\kappa_3]$ are two distinct $4$-leaf shapes in $\Sh(T)$, so $T$ is not minimal.
If $a=b=2$ then $T=[\kappa_2,\kappa_2]$, which is minimal.
If $a=1$, then $\Sh(L)=\{\lf\}\subseteq\Sh(R)$ and Lemma~\ref{lem:product} gives $s(T)=1+s(R)$, so $T$ is minimal iff $R$ is.
Hence the minimal shapes are $\lf,\kappa_2,\kappa_3$, the two $4$-leaf shapes $[\lf,\kappa_3]$, $[\kappa_2,\kappa_2]$, and for $n\ge 5$ exactly the two shapes $[\lf,[\lf,\dots,[\lf,y]\dots]]$ with $y\in\{[\lf,\kappa_3],[\kappa_2,\kappa_2]\}$.
\end{proof}

\section{Formal verification in Lean}\label{sec:lean}

We formalized the statement of Conjecture~\ref{conj:DR} together with the proofs of Theorems~\ref{thm:main} and \ref{thm:unique} and of the lower bound in Proposition~\ref{prop:min} in Lean~4 (version 4.33.1) with Mathlib (commit \texttt{0df444a}).
Trees are plane binary trees, shapes are classes for the equivalence relation generated by child swaps, and the prunings are defined by the recursion \eqref{eq:recursion}:
\begin{Verbatim}[fontsize=\small,frame=single]
inductive BTree where
  | leaf : BTree
  | node : BTree → BTree → BTree

inductive Iso : BTree → BTree → Prop
  | leaf : Iso leaf leaf
  | node {a b c d} : Iso a c → Iso b d → Iso (node a b) (node c d)
  | swap {a b c d} : Iso a d → Iso b c → Iso (node a b) (node c d)

def prunings : BTree → Finset BTree
  | leaf => {leaf}
  | node l r => insert leaf ((l.prunings ×ˢ r.prunings).image
                                fun p => node p.1 p.2)

noncomputable def rootConfigs (T : BTree) : ℕ :=
  Nat.card {s : Quotient shapeSetoid //
              ∃ P ∈ T.prunings, P ≠ leaf ∧ Quotient.mk _ P = s}

def drMax (n : ℕ) : ℕ :=
  3 ^ (n / 3) * (2 * (n % 3) ^ 2 + 8 * (n % 3) + 23) / 27

theorem disanto_rosenberg (n : ℕ) (hn : 2 ≤ n) :
    IsGreatest {k | ∃ T : BTree, T.leaves = n ∧ T.rootConfigs = k}
      (drMax n)
\end{Verbatim}
Here \texttt{shapeSetoid} is the equivalence closure of \texttt{Iso}.
The same file also proves
\begin{Verbatim}[fontsize=\small,frame=single]
theorem disanto_rosenberg_unique (n : ℕ) (hn : 9 ≤ n) (T : BTree)
    (hleaves : T.leaves = n) (hmax : T.rootConfigs = drMax n) :
    BTree.Iso T (BTree.extremal n)

theorem disanto_rosenberg_minimum (n : ℕ) (hn : 1 ≤ n) :
    IsLeast {k | ∃ T : BTree, T.leaves = n ∧ T.rootConfigs = k} (n - 1)
\end{Verbatim}
where \texttt{BTree.extremal n} is an explicit tree of shape $u_n$.
All three theorems are proved without \texttt{sorry} and without \texttt{native\_decide}, and depend only on the standard axioms \texttt{propext}, \texttt{Classical.choice} and \texttt{Quot.sound}.
In fact two independent formalizations were produced from the same fixed statement file (files \texttt{Main.lean} and \texttt{Independent.lean} in \nolinkurl{lean/KOS/Phylogenetics/AncestralConfigurations/} of the repository; about 900 and 1000 lines, each checked in under ten seconds once Mathlib is built); the second one proves the uniqueness statement in the form that any two maximizers with $n\ge 9$ leaves are isomorphic, together with the corresponding $\exists!$ statement for shapes.
In particular the base cases $9\le n\le 14$ of Theorem~\ref{thm:unique} are covered by kernel-checked finite computations inside Lean, so Theorem~\ref{thm:unique} does not depend on the external enumeration mentioned in Section~\ref{sec:unique}.

The formal proof follows Sections~\ref{sec:prunings}--\ref{sec:unique}, with two technical differences.
Shapes are represented by a complete numerical invariant: the code of $\lf$ is $0$ and the code of $[x,y]$ is $\pi(\min,\max)+1$, where $\pi$ is the Cantor pairing function applied to the codes of $x$ and $y$; isomorphism of trees is proved equivalent to equality of codes.
Lemma~\ref{lem:arith} is proved by reducing to $1\le a,b\le 11$, using $F(n+3)=3F(n)$ for $n\ge 9$ and the fact that the correction term $\binom{\min(a,b,3)}{2}$ does not change when $b\ge 3$ is replaced by $b+3$, followed by a finite check by \texttt{decide}.
The trusted part of the formalization consists of the displayed definitions; their agreement with the definitions of \cite{DR2019} is the content of Lemma~\ref{lem:reduction} and of the recursion~\eqref{eq:recursion}.

\section{Concluding remarks}\label{sec:remarks}

\begin{enumerate}
\item The product formula~\eqref{eq:product} makes $\cstar$ computable in time polynomial in the number of distinct pruned shapes, and may be useful for other extremal questions about non-equivalent configurations.
\item Disanto and Rosenberg also consider the \emph{total} number of non-equivalent configurations, summed over all internal nodes of the species tree.
Its maximum is not attained by $u_n$ for $n\ge 5$: by exhaustive computation the maxima for $n=2,\dots,17$ are
1, 3, 6, 10, 15, 22, 33, 46, 69, 94, 141, 193, 285, 397, 573, 817,
whereas $u_n$ attains 1, 3, 6, 9, 14, 21, 30, 45, 57, 93, 117, 159, 237, 327, 459, 663.
Since the total is at least the root count and at most $(n-1)M_n$ (each of the $n-1$ internal nodes contributes at most $M_n$, as $M_m$ is nondecreasing in $m$), its growth constant is also $3^{1/3}$ by Corollary~\ref{cor:growth}; determining the exact maximum and the extremal shapes is an open problem.
\end{enumerate}

\subsection*{Author contributions}
GPT-6-Astra identified the conjecture as an open problem suitable for this project and proposed the product formula~\eqref{eq:product} together with the reduction to Lemma~\ref{lem:arith}; it also produced the first Lean formalization.
Claude Opus~5.5 completed and checked the proofs (Lemma~\ref{lem:arith}, the uniqueness and minimum results), carried out the computations, wrote the manuscript and produced the second, independent Lean formalization.
Claude Fable~5.1 refereed the proof against the definitions of \cite{DR2019} and searched the literature for prior work.
A human operator initiated the project, suggested looking for open problems in applied fields, and provided computing resources; the mathematics, the formal proofs and the text were produced by the models.
All statements were checked by exhaustive computation for small $n$, and Theorems~\ref{thm:main} and \ref{thm:unique} were verified in Lean by two independent formalizations.

\begin{thebibliography}{99}

\bibitem{DR2017}
F.~Disanto and N.~A. Rosenberg,
\emph{Enumeration of ancestral configurations for matching gene trees and species trees},
J. Comput. Biol. \textbf{24} (2017), 831--850.

\bibitem{DR2019}
F.~Disanto and N.~A. Rosenberg,
\emph{On the number of non-equivalent ancestral configurations for matching gene trees and species trees},
Bull. Math. Biol. \textbf{81} (2019), 384--407. doi:10.1007/s11538-017-0342-x.

\bibitem{DFPR2022}
F.~Disanto, M.~Fuchs, A.~R. Paningbatan and N.~A. Rosenberg,
\emph{The distributions under two species-tree models of the number of root ancestral configurations for matching gene trees and species trees},
Ann. Appl. Probab. \textbf{32} (2022), 4426--4458.

\bibitem{DFHPR2024}
F.~Disanto, M.~Fuchs, C.-Y. Huang, A.~R. Paningbatan and N.~A. Rosenberg,
\emph{The distributions under two species-tree models of the total number of ancestral configurations for matching gene trees and species trees},
Adv. in Appl. Math. \textbf{152} (2024), 102594.

\bibitem{LR2024}
E.~Lappo and N.~A. Rosenberg,
\emph{A lattice structure for ancestral configurations arising from the relationship between gene trees and species trees},
Discrete Appl. Math. \textbf{343} (2024), 65--81.

\bibitem{Wu2012}
Y.~Wu,
\emph{Coalescent-based species tree inference from gene tree topologies under incomplete lineage sorting by maximum likelihood},
Evolution \textbf{66} (2012), 763--775.

\bibitem{Lean4}
L.~de Moura and S.~Ullrich,
\emph{The Lean 4 theorem prover and programming language},
in: Automated Deduction -- CADE 28, Lecture Notes in Comput. Sci. \textbf{12699}, Springer, 2021, 625--635.

\bibitem{Mathlib}
The mathlib Community,
\emph{The Lean mathematical library},
in: Proceedings of the 9th ACM SIGPLAN International Conference on Certified Programs and Proofs (CPP 2020), 367--381.

\end{thebibliography}

\end{document}
